\documentclass[10pt,a4paper]{amsart}
\usepackage[margin=19mm]{geometry}
\usepackage{amsmath,amssymb,mathtools}
\usepackage[scr=rsfs,cal=euler]{mathalfa}
\usepackage{xfrac}
\usepackage[unicode,hidelinks,bookmarks=false]{hyperref}
\newcommand{\ie}{i.e.\ }
\newcommand{\cf}{cf.\ }
\newcommand{\resp}{respectively\ }
\newcommand{\Subsection}{\subsection}
\providecommand{\nobreakdash}{\nobreak\hbox{-}}
\setlength{\emergencystretch}{2em}
\multlinegap0pt
\usepackage[shortlabels]{enumitem}
\usepackage{amssymb}
\usepackage{mathtools}
\usepackage{xcolor}
\usepackage{tikz}
\usepackage{float}
\newtheorem{lemma}{Lemma}
\usepackage{cleveref}
\crefname{lemma}{Lemma}{Lemmas}
\newcommand{\Max}{\operatorname{Max}}
\newcommand{\bbT}{\mathbb T}
\newcommand{\lessdotp}{\mathrel{\lessdot_p}}
\title{An Equivariant Landweber Exact Functor Theorem for Abelian Compact Lie Groups}
\author{Yingxin Li}
\date{Source: arXiv:2609.17519v1 (CC BY 4.0)}
\begin{document}
\maketitle

\noindent\emph{Source and licence.} An Equivariant Landweber Exact Functor Theorem for Abelian Compact Lie Groups. Version arXiv:2609.17519v1 (CC BY 4.0), distributed by arXiv under the Creative Commons Attribution 4.0 International licence, http://creativecommons.org/licenses/by/4.0/, as recorded in the arXiv licence metadata for this exact version and shown on the abstract page https://arxiv.org/abs/2609.17519v1. Full licence notice: \texttt{licenses/arxiv-2609.17519-CC-BY-4.0.txt}.\par\smallskip

\noindent\textit{Editorial scope.} Counted unit: the article's lemma condition (source0.tex lines 981--987). The article's complete original proof of it is delivered with it (source0.tex lines 988--1038, one \texttt{proof} environment of 51 lines). No mathematics is written, repaired or re-derived: the statement and the proof are the article's own lines, in the article's own notation, and no retained line is substituted or re-flowed.  Retained uncounted as the source-local context the proof needs: lines 960--966.  Nothing was omitted from any retained span, so the package carries no omission record.  No retained line contains a citation, so the wrapper carries no bibliography and no bibliography entry.  No retained line contains a figure environment, an image-inclusion command or drawing code, so no image is delivered and the package declares no asset dependency.  Editorial formatting only, in addition: an AMS \texttt{amsart} preamble that restates the article's own declarations and macros used by the retained lines, namely the environments \texttt{lemma} on separate counters, together with the standard packages those lines need.  The retained environments print in this document as Lemma 1 in order of appearance, whereas the article prints the same items with the numbering of its own full document.  The complete native source of the article as distributed in the arXiv e-print for this exact version is bundled unchanged as \texttt{sources/210612/source0.tex}.
\begin{lemma}\label{lem:compact-exact-regular}
Suppose that $(e,d)$ is an exact pair on a $D$-module $M$.
\begin{enumerate}[label=\textup{(\roman*)},leftmargin=3em]
\item If $f$ is injective on $M$, then $(e,d)$ is exact on $M/fM$.
\item Every $M$-regular sequence is regular on $M/eM$.
\end{enumerate}
\end{lemma}

\begin{lemma}\label{lem:compact exactness condition}
Let $N$ be an $L_G$-module satisfying $(\mathrm{RE})$ and $(\mathrm{EU})$. 
\begin{enumerate}[label=\textup{(\roman*)},leftmargin=3em]
  \item For every prime $p$, every closed subgroup $H\leq G$, and every $K\in\Max(H)$, the pair $(e_{H/K},d_{H/K})$ is exact on $N_{H,(p)}$. 
  \item For every closed subgroup $H\leq G$, whenever $\lambda_1,\lambda_2, \dots, \lambda_r\in H^\vee$ are linearly independent in $H^\vee\otimes \mathbb Q$, there Euler classes form a regular sequence on $N_H$.
\end{enumerate}
\end{lemma}

\begin{proof}
  (i) For the first statement, it suffices to consider the case $H/K\cong C_p$. let $A_H$ and $A_K$ be the projection of $H$ and $K$ to $A$, respectively. 
  
  (a) If $A_H=A_K=0$, then $K< H\leq T\cong \bbT^m$, and $H$ must be finite. Let 
  \[
  \Lambda_H:=\ker (T^\vee \to H^\vee), \Lambda_K:=\ker ( T^\vee\to K^\vee).
  \]
  We can choose a generator of $\Lambda_H$ given by $\{\lambda_1,\dots, \lambda_r\}$ such that the $r$-tuple $(\lambda_1,\dots,\lambda_r)$ are linearly independent, and $\Lambda_K=\langle p\lambda_1,\lambda_2,\dots, \lambda_r\rangle$. Then
  \[
  N_H \cong N_T/(e_{p\lambda_1},e_{\lambda_2}, \dots, e_{\lambda_r}),\quad N_K \cong N_T/(e_{\lambda_1},e_{\lambda_2}, \dots, e_{\lambda_r})
  \]
  and $e_{H/K}=\operatorname{Res}_H^T e_{\lambda_1}$. Since $e_{p\lambda_1}$ restrict to $0$ in $L_H$, $e_{p\lambda_1}=e_{\lambda_1}\cdot \psi_p$ such that $d_{H/K}=\operatorname{Res}_H^T \psi_p$. Since $e_{\lambda_1}$ is a nonzero divisor in $N_T$. Thus $(e_{\lambda_1}, \psi_p)$ is an exact pair of zero divisor in $N_T/e_{p\lambda_1}$. Note that $(\mathrm{RE})$ implies that $(e_{\lambda_2},\dots, e_{\lambda_r})$ is regular on $N_T/(e_{p\lambda_1})$, part~\textup{(i)} of \Cref{lem:compact-exact-regular} implies that $(e_H/K,d_{H/K})$ is exact on $N_H$. 

  (b) If $A_H=A_K\neq 0$, then $H\cong L\times B$ and $K\cong  W\times B$ for some $W\lessdotp L< T$ and $B\leq A$. For any subgroup $B\leq A$, choose a chain
  \[
  B=B_0<B_1<\cdots<B_r=A
  \]
  in which each $B_i$ is a maximal subgroup of $B_{i+1}$, so that $B_{i+1}/B_i\cong C_{q_i}$ for some prime $q_i$. At each step, restriction from $T\times B_{i+1}$ to $T\times B_i$ kills the Euler class $e_{B_{i+1}/B_i}\in L_{T\times B_{i+1}}$, which belongs to an exact pair. Applying part~\textup{(ii)} of \Cref{lem:compact-exact-regular} successively, we may therefore assume without loss of generality that $B=A$. Let 
  $$\Lambda_{K}:=\ker (G^\vee \to {K}^\vee)=\langle \lambda_1,\lambda_2,\dots,\lambda_r\rangle, \quad \Lambda_{H}=\Lambda_{K}/p\lambda_1.$$ 
  Here $\lambda_i=\tau_i\alpha_i$ for some $\tau_i\in T^\vee$ and $\alpha_i\in A^\vee$. 
  Since $A$ is finite, we can choose $s\in \mathbb Z$ such that $s A^\vee=0$. Then $$e_{s\lambda_i}=e_{s\tau_i}.$$
  Condition $(\mathrm{RE})$ ensures that the sequence $(e_{s\tau_1},\dots,e_{s\tau_r})$ is regular on $N$. For each $i$, we have $e_{s\lambda_i}=e_{\lambda_i}u_i$ for some $u_i\in L_G$, and hence the sequence $(e_{\lambda_1},\dots,e_{\lambda_r})$ is also regular on $N$. As in the case $K<H\leq T\cong\bbT^m$, it then follows from part~\textup{(i)} of \Cref{lem:compact-exact-regular} that $(e_{H/K},d_{H/K})$ is an exact pair on $N_H$.

  (c) Now assume that $A_H/A_K\cong C_p$. If $H\cong \bbT^r\times A$ and $K\cong \bbT^r\times B$ for some $B\lessdotp A$, then $(e_{H/K},d_{H/K})$ is exact on $N_H$, which is followed by $(\mathrm{EU})$ and \Cref{lem:compact-exact-regular}. In general, let $C_H=T\times A_H$ and $C_K=T\times A_K$. Then the pullback of $e_{A_H/A_K}$ belongs to an exact pair in $N_{C_H}$. As the case (b) above, we can choose a regular sequence of Eulers $(e_1,\dots,e_r)$ such that 
  \[
  N_H\cong N_{C_H}/(e_1,\dots, e_r).
  \]
  Since $K=H\cap C_K$, the quotient character $C_H \to A_H/A_K\cong C_p$ restricts to the quotient $H\to H/K$, thus $e_{H/K}$ is the restriction of $e_{C_H/C_K}$ and thus $(e_{H/K},d_{H/K})$ is exact on $N_H$ by part (i) of \Cref{lem:compact-exact-regular}.

  (ii) Let $A_H$ be the projection of $H$ to $A$ and let $C=T\times A_H$. Let 
  \[
  \Lambda_H =\ker (C^\vee\to H^\vee).
  \]
  $\Lambda_H$ is free since an element in $\Lambda_H$ belongs to \(A_H^\vee\) and vanishes on $H$ must be trivial since $H\to A_H$ is surjective. Then we can choose a basis $\chi_1,\dots, \chi_s$ such that $e_i=e_{\chi_i}$ form a regular sequence on $N_C$, and 
  \[
  N_H \cong N_C/(e_1,\dots, e_s).
  \]
  Suppose
  \[
  \lambda_1,\ldots,\lambda_r\in H^\vee
  \]
  are linearly independent in \(H^\vee\otimes\mathbb Q\), and choose lifts
  \[
  \widetilde\lambda_1,\ldots,\widetilde\lambda_r\in C^\vee.
  \]
  Then
  \[
  \widetilde\lambda_1,\ldots,\widetilde\lambda_r, \chi_1,\ldots,\chi_s
  \]
  are linearly independent in $ C^\vee\otimes\mathbb Q\cong T^\vee\otimes\mathbb Q$, and the sequence of Euler classes associated to these characters is regular on $N_C$. Then $e_{\lambda_1}, \dots, e_{\lambda_r}$ is regular on $N_H$. 
\end{proof}

\end{document}
